\documentclass[aspectratio=169]{beamer}
\usepackage[T1]{fontenc}
\usepackage{amsmath}
\usepackage{graphicx}
\usetheme{default}
\setbeamertemplate{navigation symbols}{}
\newcommand{\E}{\mathbb{E}}
\newcommand{\norm}[1]{\left\lVert #1\right\rVert}
\begin{document}
\title{Presenting with OverLyX}
\subtitle{An ordinary beamer deck, edited as you see it}
\author{@@NAME@@}
\institute{Made with OverLyX}
\section{Getting started}

\begin{frame}{This is an ordinary beamer file}
\begin{itemize}
\item The file is plain \LaTeX{} (\texttt{slides.tex}): it compiles anywhere, and desktop \LyX{} opens it too.
\item Every frame is a paragraph of the layout \emph{Frame}; its title is the argument in the frame's first line.
\item Edit the text right here --- lists, formulas and figures look as they will in the PDF.
\item Compile the PDF with the build button; the PDF panel shows the result.
\end{itemize}
\note[item]{Speaker notes are written with \texttt{\textbackslash note} in a frame: they show in the presenter view (press S while presenting), not on the slide.}
\end{frame}

\begin{frame}{Presenting}
\begin{itemize}
\item Press \textbf{F5} to present from the first slide (\textbf{Shift+F5}: from the current one).
\item \textbf{Space}, \textbf{Right} or a click go forward, \textbf{Left} goes back, \textbf{Esc} ends.
\item \textbf{S} opens the presenter view: this slide, the next one, your notes and a timer.
\item \textbf{B} / \textbf{W} blank the screen, \textbf{L} is a laser pointer, a number + \textbf{Enter} jumps there.
\end{itemize}
\note[item]{Try it now: press F5.}
\end{frame}

\begin{frame}{Or a canvas deck}
This deck is linear: beamer lays out the text of each frame. OverLyX also makes \emph{canvas decks}, as in PowerPoint or Keynote, where every title, text box, shape and picture is an object you place on the slide yourself.
\begin{itemize}
\item Open the project \emph{Example: slide deck} to see one --- its speaker notes are a guided tour.
\item Start your own with File $\triangleright$ New slides / poster / page\ldots{}
\end{itemize}
\end{frame}

\section{Overlays}

\begin{frame}{One item at a time}
\begin{itemize}[<+->]
\item \texttt{[<+->]} after \texttt{\textbackslash{}begin\{itemize\}} uncovers the items one by one.
\item The presentation shows the same steps as the PDF, which gets one page per step.
\item Formulas work in every step: $\E[X]=\sum_{x}x\,p(x)$.
\end{itemize}
\end{frame}

\begin{frame}{Pauses and item specifications}
\begin{itemize}
\item<1-> \texttt{\textbackslash{}item<1->} is there from the first step.
\item<2-> \texttt{\textbackslash{}item<2->} appears with the second.
\item<3-> \texttt{\textbackslash{}item<3->} with the third.
\end{itemize}
A \texttt{\textbackslash{}pause} stops here \pause and the rest follows on the next step.
\end{frame}

\begin{frame}{Only, uncover, alert and alt}
\only<1>{\texttt{\textbackslash{}only<1>} --- this sentence is there on the first step only.}\only<2->{\texttt{\textbackslash{}only<2->} --- and this one replaces it from the second.}

\uncover<3->{\texttt{\textbackslash{}uncover<3->} keeps its room while covered and appears on step three.}

\alert<2>{\texttt{\textbackslash{}alert<2>} turns red on step two.}

\alt<4>{\textbf{On step four, alt shows its first text.}}{\texttt{\textbackslash{}alt<4>\{\ldots{}\}\{\ldots{}\}} shows its second text on the other steps.}
\end{frame}

\section{Structure}

\begin{frame}{Columns and figures}
\begin{columns}
\begin{column}{0.45\textwidth}
Columns sit side by side, as in the PDF.
\begin{itemize}
\item Their widths come from \texttt{\textbackslash{}begin\{column\}\{0.45\textbackslash{}textwidth\}}.
\item A figure is an ordinary \texttt{\textbackslash{}includegraphics}.
\end{itemize}
\end{column}
\begin{column}{0.5\textwidth}
\includegraphics[width=\linewidth]{figures/decay}
\end{column}
\end{columns}
\end{frame}

\begin{frame}{Blocks}
\begin{block}{A block}
Blocks have a title and a body.
\end{block}
\begin{alertblock}<2->{An alert block}
This one appears on the second step: blocks take overlay specifications too.
\end{alertblock}
\begin{exampleblock}<3->{An example block}
And this one on the third.
\end{exampleblock}
\end{frame}

\begin{frame}{Theorems and formulas}
\begin{theorem}[Jensen]
For a convex function $\varphi$ and an integrable random variable $X$,
\[
\varphi\left(\E[X]\right)\le\E\left[\varphi(X)\right].
\]
\end{theorem}
\begin{proof}
Take a tangent line $\ell(x)=a+bx\le\varphi(x)$ at $\E[X]$ and use the linearity of $\E$.
\end{proof}
Macros from the preamble render in the editor: $\norm{x-\E[X]}^{2}$.
\end{frame}

\begin{frame}{A table}
\begin{center}
\begin{tabular}{lcc}
\hline
Method & Steps & Error\tabularnewline
\hline
Euler & 1000 & $3.1\times10^{-2}$\tabularnewline
Heun & 500 & $4.7\times10^{-4}$\tabularnewline
Runge--Kutta 4 & 100 & $2.2\times10^{-7}$\tabularnewline
\hline
\end{tabular}
\end{center}
Tables are edited cell by cell; Tab moves to the next cell.
\end{frame}

\begin{frame}{Your turn}
\begin{enumerate}
\item Change a word in this frame and watch the PDF follow.
\item Make a new frame: put the cursor after a frame, press Enter and choose the layout \emph{Frame}.
\item Copy an item with an overlay from the frames above and change its number --- or edit the \LaTeX{} source in the source pane.
\item Press F5 and present it.
\end{enumerate}
\end{frame}

\end{document}
